\newpage
\section{IMPLEMENTATION DETAILS}

This final deployment of the suggested system consists of mechanical actuation, RSSI embedding,
offline RL, and field testing in one chain. This final process consists of four
intertwined phases: hardware design, embedded software implementation, offline
training, and field experimentations.

Hardware-wise, a mechanism of two degrees of freedom was constructed to rotate the
receiving antenna in pan and tilt while maintaining power and data communication
through the rotating part of the interface. Embedded-wise, ESP32-S3 was chosen as a
primary processor of packets received from the wireless network, the extraction of
RSSI information, as well as control over servos and steppers.

\vspace{0.5cm}
\begin{table}[H]
\centering
\caption{Implementation Stages of the Proposed System}
\label{tab:implementation_stages}
\renewcommand{\arraystretch}{1.4}
\begin{tabular}{|
>{\raggedright\arraybackslash}p{3.2cm}|
>{\raggedright\arraybackslash}p{4.2cm}|
>{\raggedright\arraybackslash}p{5.2cm}|
}
\hline
\textbf{Stage} & \textbf{Primary Output} & \textbf{Purpose} \\
\hline
Mechanical and Electronic Circuit Implementation & Two-DOF antenna platform, slip-ring interface, regulated power distribution & To realize stable pan--tilt motion and support practical RSSI measurements \\
\hline
Firmware Implementation & ESP32-S3 control loop, RSSI averaging, motor actuation, experiment logging & To execute the alignment algorithms under embedded constraints \\
\hline
Offline RL Training & Trained Q-table and tuned hyperparameters & To learn a compact policy that can be deployed without online training \\
\hline
Practical Validation & Indoor and outdoor trial records at 10~m to 30~m & To generate consistent experimental data for comparing the deployed RL controller with the baseline methods \\
\hline
\end{tabular}
\end{table}

The implementation details are mentioned in following subsections describe how these stages were
realized and how final prototype was prepared for experimental evaluation.

